\documentclass[11pt]{article}
\usepackage{amsmath}
\usepackage{amssymb}
\usepackage{geometry}
\geometry{margin=1in}

\title{Plan 02-02: Gauss vs. Lagrange Formulation}
\author{J2+J3+Drag Secular Perturbation Theory}
\date{2026-04-02 — Phase 2}

\begin{document}
\maketitle

\section{Objective}
Decide whether to use Gauss planetary equations or equivalent disturbing potential for drag perturbations.

\section{Lagrange Planetary Equations (Conservative Forces)}

For conservative perturbations (gravitational harmonics), we used:
\begin{equation}
\frac{dx}{dt} = f(R, \text{elements})
\end{equation}
where $R$ is a disturbing potential function.

**Requirement:** $R$ must exist as a scalar potential (conservative force).

\section{Gauss Planetary Equations (Non-Conservative Forces)}

For non-conservative forces (drag, thrust), use Gauss form:
\begin{align}
\frac{da}{dt} &= \frac{2a^2}{h}\left[e\sin f\, a_R + \frac{p}{r} a_T\right]\\
\frac{de}{dt} &= \frac{1}{h}\left[\sin f\, a_R + ((\cos f + \cos E) a_T\right]
\end{align}

where:
\begin{itemize}
    \item $a_R$ — radial acceleration component
    \item $a_T$ — tangential acceleration component (along velocity)
    \item $h = \sqrt{\mu p} = \sqrt{\mu a(1-e^2)}$ — specific angular momentum
    \item $p = a(1-e^2)$ — semi-latus rectum
\end{itemize}

\section{Can Drag Have an Equivalent Potential?}

**Analysis:** Drag force is $\vec{F}_D = -\frac{1}{2}\rho v^2 B\,\hat{v}$ (velocity-dependent).

For a potential $R_{\text{drag}}$ to exist:
\begin{equation}
\vec{F}_D = -\nabla R_{\text{drag}}
\end{equation}

must hold. But drag depends on $\vec{v}$, not just position $\vec{r}$.

**Conclusion:** \textbf{No true scalar potential exists for drag} (velocity-dependent, dissipative).

\section{Decision: Use Gauss Equations}

\textbf{Chosen formulation:} Gauss planetary equations

\textbf{Rationale:}
\begin{enumerate}
    \item Drag is dissipative → no conservative potential
    \item Gauss equations explicitly handle non-conservative forces
    \item Well-established in literature (King-Hele 1987, Vallado 2013)
    \item Allows direct treatment of tangential acceleration
\end{enumerate}

\section{Simplified Form for Circular Orbits}

For near-circular sun-synchronous orbits ($e \ll 1$):
\begin{itemize}
    \item $a_R \approx 0$ (drag is tangential)
    \item $a_N = 0$ (symmetric drag, no out-of-plane component)
    \item $a_T \approx -\frac{1}{2}\rho v^2 B$ (opposes velocity)
\end{itemize}

Gauss equations simplify to:
\begin{align}
\frac{da}{dt} &\approx \frac{2a^2}{h} \cdot \frac{p}{r} a_T = \frac{2a}{\sqrt{\mu a}} a_T = \frac{2a^{3/2}}{\sqrt{\mu}} a_T\\
\frac{de}{dt} &\approx \frac{1}{h}(1 + \cos E) a_T\\
\frac{di}{dt} &= 0 \quad \text{(no normal component)}
\end{align}

\section{Success Criteria}
\begin{itemize}
    \item[$\checkmark$] Both formulations analyzed
    \item[$\checkmark$] Decision made: \textbf{Gauss equations}
    \item[$\checkmark$] Rationale documented
    \item[$\checkmark$] Simplified circular-orbit form identified
\end{itemize}

\end{document}
